% Baccalauréat malgache, série D, 2016 : corrigé.
\TCDBacAttribution{D}{2016}
\providecommand{\Courbe}{(\mathscr{C})}
\providecommand{\CourbeP}{(\mathscr{C}')}
\providecommand{\D}{(\mathscr{D})}
\providecommand{\cm}{~\mathrm{cm}}
\providecommand{\dx}{~\mathrm{d}x}
%
%%%%%%%%%%%%%%%%%%%%%%%%%%%%%%%%%%%%%%%%%%%%%%%%%%%%%%%%%%%%%%%%%%%%%%%%%%%%%%%%%%%%%%%%%%%%%%%
\noindent \underline{\textsc{\textbf{Exercice I}}} (5 points)
	\begin{enumerate}
		\item \begin{enumerate}
			\item \vln{Calculs :}
				$P(1) = 1 - 2 + 2 + i - 1 - i$ et $(1-2i)^2 = 1 - 4i + 4i^2$.

				\ansbox{$P(1) = 0$ et $(1-2i)^2 = -3 - 4i$.}
			\item \vln{Résolution :}
				$1$ est racine, donc $P(z) = (z-1)(z^2 + pz + q) = z^3 + (p-1)z^2 + (q-p)z - q$.
				Par identification : $p - 1 = -2$, $-q = -1 - i$, d'où $p = -1$, $q = 1 + i$ ; vérification : $q - p = 2 + i$ \checkmark.

				$z^2 - z + 1 + i = 0$ a pour discriminant $\Delta = 1 - 4(1+i) = -3 - 4i = (1-2i)^2$, d'où
				$z = \frac{1 + (1-2i)}{2} = 1 - i$ et $z = \frac{1 - (1-2i)}{2} = i$.

				\ansbox{$S_{\mathbb{C}} = \{1\ ;\ 1 - i\ ;\ i\}$.}
		\end{enumerate}
		\item \vln{Forme trigonométrique de $a$ :}
			$a = \frac{1-i}{i} = \frac{(1-i)(-i)}{1} = -i + i^2 = -1 - i$, de module $\sqrt{2}$ et d'argument $-\frac{3\pi}{4}$.

			\ansbox{$a = \sqrt{2}\left(\cos\left(-\frac{3\pi}{4}\right) + i\sin\left(-\frac{3\pi}{4}\right)\right)$.}

			D'après la formule de Moivre, $a^4 = \left(\sqrt{2}\right)^4\left(\cos(-3\pi) + i\sin(-3\pi)\right) = 4\times(-1)$.

			\ansbox{$a^4 = -4$.}
		\item \begin{enumerate}
			\item \vln{Éléments caractéristiques de $S$ :}
				$1 - i = \sqrt{2}\left(\cos\left(-\frac{\pi}{4}\right) + i\sin\left(-\frac{\pi}{4}\right)\right)$ : rapport $\sqrt{2}$, angle $-\frac{\pi}{4}$.
				Le centre vérifie $\omega = (1-i)\omega + 2i$, soit $i\omega = 2i$, $\omega = 2$.

				\ansbox{$S$ : centre $\Omega(2)$, rapport $\sqrt{2}$, angle $-\frac{\pi}{4}$.}
			\item \vln{Nature de $S^4$ :}
				$S$ s'écrit $z' - 2 = (1-i)(z-2)$, donc $S^4$ s'écrit $z' - 2 = (1-i)^4(z-2)$.
				Or $(1-i)^2 = -2i$ et $(1-i)^4 = (-2i)^2 = -4$. Ainsi $S^4 : z' = -4z + 10$.

				\ansbox{$S^4$ est l'homothétie de centre $\Omega(2)$ et de rapport $-4$}

				(c'est aussi la similitude directe de centre $\Omega$, de rapport $4 = \left(\sqrt{2}\right)^4$ et d'angle $4\times\left(-\frac{\pi}{4}\right) = -\pi$).
			\item \vln{Construction :}
				$z_{A'} = -4 + 10 = 6$, $z_{B'} = -4(1-i) + 10 = 6 + 4i$, $z_{C'} = -4i + 10 = 10 - 4i$.

				\ansbox{$A'(6;0)$, $B'(6;4)$, $C'(10;-4)$.}

				(Figure à l'échelle $\frac{1}{4}$ : une unité est représentée par $0{,}5\cm$ au lieu de $2\cm$.)
				\begin{center}
				\begin{tikzpicture}[scale=.5,>=stealth]
					\draw[very thin,gray!40] (-1,-5) grid (11,5);
					\draw[->] (-1,0) -- (11.3,0) node[below] {$x$};
					\draw[->] (0,-5) -- (0,5.3) node[left] {$y$};
					\draw[thick,loko1] (1,0) -- (1,-1) -- (0,1) -- cycle;
					\draw[thick,loko3] (6,0) -- (6,4) -- (10,-4) -- cycle;
					\draw[gray,dashed] (1,0) -- (6,0) (1,-1) -- (6,4) (0,1) -- (10,-4);
					\foreach \p/\n/\pos in {(1,0)/A/above left,(1,-1)/B/below,(0,1)/C/above right,(6,0)/A'/below right,(6,4)/B'/above,(10,-4)/C'/below,(2,0)/\Omega/below}
						\fill \p circle (3pt) node[\pos] {$\n$};
				\end{tikzpicture}
				\end{center}
		\end{enumerate}
	\end{enumerate}
%%%%%%%%%%%%%%%%%%%%%%%%%%%%%%%%%%%%%%%%%%%%%%%%%%%%%%%%%%%%%%%%%%%%%%%%%%%%%%%%%%%%%%%%%%%%%%%
\noindent \underline{\textsc{\textbf{Exercice II}}} (5 points)
	\begin{enumerate}
		\item \vln{Probabilités des numéros :}
			les $6$ faces sont équiprobables (probabilité $\frac{1}{6}$ chacune) ; le numéro 2 est porté par deux faces.
			(Le dé est << truqué >> par sa numérotation, non par ses faces.)

			\ansbox{$P_1 = P_3 = P_4 = P_6 = \frac{1}{6}$ et $P_2 = \frac{2}{6} = \frac{1}{3}$.}
		\item Les deux lancers sont indépendants : la probabilité d'un couple $(i;j)$ de numéros est $P_i \times P_j$.
			\begin{enumerate}
				\item \vln{Événement $A$ :}
					$4 = 1 + 3 = 3 + 1 = 2 + 2$, donc
					$P(A) = P_1P_3 + P_3P_1 + P_2P_2 = \frac{1}{36} + \frac{1}{36} + \frac{1}{9} = \frac{6}{36}$.

					\ansbox{$P(A) = \frac{1}{6}$.}
				\item \begin{itemize}
					\item \vln{Loi de $X$ :}
						en comptant les $36$ couples de faces (le 2 compte pour deux faces) :
						$X = 2$ : $(1;1)$, $1$ cas ; $X = 3$ : $(1;2)$, $(2;1)$, $4$ cas ; $X = 4$ : $6$ cas ;
						$X = 5$ : $(1;4)$, $(4;1)$, $(2;3)$, $(3;2)$, $2 + 4 = 6$ cas ; $X = 6$ : $(2;4)$, $(4;2)$, $(3;3)$, $4 + 1 = 5$ cas ;
						$X = 7$ : $(1;6)$, $(6;1)$, $(3;4)$, $(4;3)$, $4$ cas ; $X = 8$ : $(2;6)$, $(6;2)$, $(4;4)$, $5$ cas ;
						$X = 9$ : $(3;6)$, $(6;3)$, $2$ cas ; $X = 10$ : $(4;6)$, $(6;4)$, $2$ cas ; $X = 12$ : $(6;6)$, $1$ cas
						($X = 11$ est impossible).
						\[
						\begin{array}{|c|c|c|c|c|c|c|c|c|c|c|}
							\hline
							x_i & 2 & 3 & 4 & 5 & 6 & 7 & 8 & 9 & 10 & 12 \\ \hline
							P(X=x_i) & \frac{1}{36} & \frac{1}{9} & \frac{1}{6} & \frac{1}{6} & \frac{5}{36} & \frac{1}{9} & \frac{5}{36} & \frac{1}{18} & \frac{1}{18} & \frac{1}{36} \\[.5ex] \hline
						\end{array}
						\]
						(Vérification : $1 + 4 + 6 + 6 + 5 + 4 + 5 + 2 + 2 + 1 = 36$.)
					\item \vln{Espérance :}
						$E(X) = \frac{2 + 12 + 24 + 30 + 30 + 28 + 40 + 18 + 20 + 12}{36} = \frac{216}{36}$.

						\ansbox{$E(X) = 6$}

						(ce qui est bien le double de l'espérance d'un lancer : $\frac{1+2+2+3+4+6}{6} = 3$).
				\end{itemize}
				\item On répète $3$ fois, de façon indépendante, une épreuve dont le << succès >> (obtenir 2) a la probabilité $\frac{1}{3}$ :
					$Y$ suit la loi binomiale $\mathscr{B}\left(3;\frac{1}{3}\right)$.
					\begin{enumerate}
						\item \vln{Loi de $Y$ :} $P(Y = k) = \Comb{3}{k}\left(\frac{1}{3}\right)^k\left(\frac{2}{3}\right)^{3-k}$.
							\[
							\begin{array}{|c|c|c|c|c|}
								\hline
								k & 0 & 1 & 2 & 3 \\ \hline
								P(Y=k) & \frac{8}{27} & \frac{4}{9} & \frac{2}{9} & \frac{1}{27} \\[.5ex] \hline
							\end{array}
							\]
						\item \vln{Variance :} $V(Y) = np(1-p) = 3\times\frac{1}{3}\times\frac{2}{3}$.

							\ansbox{$V(Y) = \frac{2}{3}$.}
					\end{enumerate}
			\end{enumerate}
	\end{enumerate}
%%%%%%%%%%%%%%%%%%%%%%%%%%%%%%%%%%%%%%%%%%%%%%%%%%%%%%%%%%%%%%%%%%%%%%%%%%%%%%%%%%%%%%%%%%%%%%%
\noindent \underline{\textsc{\textbf{Probleme}}} (10 points)
	\begin{description}
		\item[A.]
			\begin{enumerate}
				\item \begin{enumerate}
					\item \vln{Sens de variation de $g$ :}
						$g'(x) = -\expo{-x} - (1-x)\expo{-x} = (x-2)\expo{-x}$, du signe de $x - 2$.

						\ansbox{$g$ est décroissante sur $]-\infty;2]$ et croissante sur $[2;+\infty[$.}

						Son minimum est $g(2) = 1 - \expo{-2} \simeq 0{,}86$.
					\item \vln{Signe de $g$ :} le minimum $1 - \expo{-2}$ est strictement positif ($\expo{-2} < 1$).

						\ansbox{$g(x) > 0$ pour tout réel $x$.}
				\end{enumerate}
				\item \vln{Variations de $f$ :}
					$f'(x) = \expo{-x} - x\expo{-x} + 1 = (1-x)\expo{-x} + 1 = g(x) > 0$ : $f$ est strictement croissante sur $\mathbb{R}$.
					Limites : en $-\infty$, $x\expo{-x} \to -\infty$ (produit de $x \to -\infty$ par $\expo{-x} \to +\infty$), donc $f(x) \to -\infty$ ;
					en $+\infty$, $x\expo{-x} \to 0$ donc $f(x) \to +\infty$.
					\[
					\begin{array}{c|ccc}
						x & -\infty & & +\infty \\ \hline
						f'(x) & & + & \\ \hline
						f(x) & -\infty & \nearrow & +\infty
					\end{array}
					\]
				\item \begin{enumerate}
					\item \vln{Branche infinie en $-\infty$ :}
						$\frac{f(x)}{x} = \expo{-x} + 1 - \frac{1}{x} \to +\infty$ quand $x \to -\infty$.

						\ansbox{$\Courbe$ admet en $-\infty$ une branche parabolique de direction $(Oy)$.}
					\item \vln{Asymptote en $+\infty$ :}
						$f(x) - (x-1) = x\expo{-x} \to 0$ quand $x \to +\infty$.

						\ansbox{$\D : y = x - 1$ est asymptote oblique à $\Courbe$ en $+\infty$.}

						\vln{Position :} $f(x) - (x-1) = x\expo{-x}$ a le signe de $x$ :
						$\Courbe$ est au-dessous de $\D$ sur $]-\infty;0[$, au-dessus sur $]0;+\infty[$, et les coupe au point $(0;-1)$.
				\end{enumerate}
				\item \vln{Tangente parallèle à $\D$ :}
					il faut $f'(x) = 1$, soit $(1-x)\expo{-x} = 0$, c'est-à-dire $x = 1$ (seule solution, car $\expo{-x} \neq 0$).
					$f(1) = \expo{-1}$.

					\ansbox{$A\left(1\ ;\ \frac{1}{\expo{}}\right)$ et $(T) : y = x - 1 + \frac{1}{\expo{}}$.}
				\item \vln{Construction :} (figure à l'échelle $\frac{1}{2}$ ; la courbe $\CourbeP$ de la question 7.c y figure aussi).
					\begin{center}
					\begin{tikzpicture}[x=1cm,y=1cm,>=stealth]
						\draw[very thin,gray!40] (-5,-5) grid (5,5);
						\draw[->] (-5,0) -- (5.3,0) node[below] {$x$};
						\draw[->] (0,-5) -- (0,5.3) node[left] {$y$};
						\draw (1,0) node[below] {\small $1$} (0,1) node[left] {\small $1$};
						\clip (-5,-5) rectangle (5,5);
						\draw[gray,dashed] (-5,-5) -- (5,5);
						\draw[loko3,thick] plot[domain=-4:5] (\x,{\x-1});
						\draw[gray,thick] plot[domain=-4:5] (\x,{\x-1+exp(-1)});
						\draw[loko1,very thick,smooth,samples=100] plot[domain=-1.06:5] (\x,{\x*exp(-\x)+\x-1});
						\draw[loko5,very thick,smooth,samples=100] plot[domain=-1.06:5] ({\x*exp(-\x)+\x-1},\x);
						\fill (1,0.368) circle (2pt) node[above left] {$A$};
						\node[loko3] at (4.6,3.2) {$\D$};
						\node[gray] at (3.3,3.4) {$(T)$};
						\node[loko1] at (-0.4,-4.5) {$\Courbe$};
						\node[loko5] at (-4.2,-0.5) {$\CourbeP$};
					\end{tikzpicture}
					\end{center}
				\item \vln{Aire :}
					sur $[0;1]$, $\Courbe$ est au-dessus de $\D$ et $f(x) - (x-1) = x\expo{-x}$.
					Par parties ($u = x$, $v' = \expo{-x}$) :
					\[
						\int_0^1 x\expo{-x}\dx = \Big[-x\expo{-x}\Big]_0^1 + \int_0^1\expo{-x}\dx = -\expo{-1} + 1 - \expo{-1} = 1 - \frac{2}{\expo{}}.
					\]
					L'unité d'aire vaut $2\cm\times2\cm = 4\cm^2$.

					\ansbox{$\mathscr{A} = 4\left(1 - \frac{2}{\expo{}}\right)\cm^2 = \left(4 - \frac{8}{\expo{}}\right)\cm^2 \simeq 1{,}06\cm^2$.}
				\item \begin{enumerate}
					\item \vln{Bijection :} $f$ est continue et strictement croissante sur $\mathbb{R}$, de limites $-\infty$ et $+\infty$.

						\ansbox{$f$ réalise une bijection de $\mathbb{R}$ sur $J = \mathbb{R}$.}
					\item \vln{Calculs :} $f(1) = \expo{-1} + 1 - 1$.

						\ansbox{$f(1) = \frac{1}{\expo{}}$.}

						$f$ est dérivable en $1$ avec $f'(1) = g(1) = 1 \neq 0$, donc
						$\left(f^{-1}\right)'\left(\frac{1}{\expo{}}\right) = \frac{1}{f'\left(f^{-1}\left(\frac{1}{\expo{}}\right)\right)} = \frac{1}{f'(1)}$.

						\ansbox{$\left(f^{-1}\right)'\left(\frac{1}{\expo{}}\right) = 1$.}
					\item \vln{Courbe $\CourbeP$ :} symétrique de $\Courbe$ par rapport à la droite $y = x$ (voir la figure) ;
						elle passe par $\left(\frac{1}{\expo{}};1\right) \simeq (0{,}4;1)$ où sa tangente, symétrique de $(T)$, a pour coefficient directeur $1$.
				\end{enumerate}
			\end{enumerate}
		\item[B.]
			\begin{enumerate}
				\item \vln{Calcul de $U_n$ :}
					$\int_0^n\expo{-x}\dx = \Big[-\expo{-x}\Big]_0^n = 1 - \expo{-n}$, donc $U_n = 1 - \left(1 - \expo{-n}\right)$.

					\ansbox{$U_n = \expo{-n}$ pour tout $n \in \mathbb{N}$} (ce qui redonne bien $U_0 = 1$).
				\item \vln{Suite géométrique :}
					pour tout $n$, $U_{n+1} = \expo{-(n+1)} = \expo{-1}\times\expo{-n} = \expo{-1}U_n$.

					\ansbox{$(U_n)$ est géométrique de raison $q = \frac{1}{\expo{}}$ et de premier terme $U_0 = 1$.}
			\end{enumerate}
	\end{description}
